\begin{ledgroupsized}[r]{120mm}
\footnotesize
\pstart
\noindent\textbf{Überlieferung:}
\pend
\end{ledgroupsized}
\begin{ledgroupsized}[r]{114mm}
\footnotesize
\pstart
\parindent -6mm
\makebox[6mm][l]{\textit{L}}%
Konzept:
LH~XXXV~10, 14~Bl.~5. 
Ein Blatt~4\textsuperscript{o};
\href{\Scan{LH_35_10_14_005wz}}%
{Wasserzeichenfragment} am Blattrand;
Rechter und unterer Rand beschnitten.
Einzweidrittel Seiten.
In der oberen linken Ecke von Bl.~5~r\textsuperscript{o}, in Rot, der Vermerk von unklarer Hand \textit{V} 
(vgl. die Träger der benachbarten Signaturen LH~XXXV~10, 14 Bl.~2 bis 13, u.\,a. 
N.\;\ref{60117},
\ref{60118},
\ref{60115},
\ref{60116},
\ref{60113},
\ref{60114},
\ref{60110},
\ref{41824}
und \ref{60111} 
in diesem Band).
\pend
\end{ledgroupsized}
\vspace*{8mm}
\selectlanguage{latin}
\pstart
\noindent
\href{\Scan{LH_35_10_14_005r}}%
{\lbrack5~r\textsuperscript{o}\rbrack}\ 
Sit Globus cujus centrum figurae\indexS{centrum figurae} \textit{C}, centrum
\edtext{gravitatis\indexS{centrum gravitatis} \textit{G}. Is tangat\indexS{globus planum inclinatum tangens}}{%
\lemma{gravitatis \textit{G}}%
\Bfootnote{%
\textit{(1)}~tangens %
\textit{(2)}~. Is tangat~\textit{L}}} 
planum\indexS{planum inclinatum}
\edtext{inclinatum \textit{AD}, in \textit{D} ita ut \textit{C}, \textit{G}, \textit{D} sint in directum, seu recta \textit{CG} continuata cadat in \textit{D}}{%
\lemma{inclinatum \textit{AD},}%
\Bfootnote{%
\textit{(1)}~puncto \textit{D} %
\textit{(2)}~in \textit{D} %
\textit{(a)}~quod sit in recta \textit{CGD} continuata %
\textit{(b)}~ita \bracks{...} in \textit{D}~\textit{L}}}. 
Ponatur autem ibi in \textit{D}, globus nonnihil truncatus,
\edtext{vel aliter impeditus}{%
\lemma{}%
\Bfootnote{vel aliter impeditus \textit{erg.~L}}} 
ita
\edtext{ut non possit volvi}{%
\lemma{ut}%
\Bfootnote{%
\textit{(1)}~globus p %
\textit{(2)}~non possit %
\textit{(a)}~volver %
\textit{(b)}~volvi~\textit{L}}} 
super plano,\indexS{planum inclinatum} sed
\edtext{descendens\indexS{globus descendens in plano sine volutione} eodem semper loco \textit{D} ipsum planum tangat\bracks{,} porro descendere nitens\indexS{globus descendere nitens}}{%
\lemma{descendens}%
\Bfootnote{%
\textit{(1)}~super ipso %
\textit{(2)}~glis %
\textit{(3)}~eodem semper %
\textit{(a)}~puncto %
\textit{(b)}~loco \textit{D} %
\textit{(aa)}~plano %
\textit{(bb)}~pla %
\textit{(cc)}~ipsum planum tangat porro
\textit{(aa)}~descendens urg %
\textit{(bb)}~descendere nitens~\textit{L}}} 
globus urgeat \textit{BF} brachium\indexS{brachium librae} librae\indexS{libra} 
\textit{EBF}, cujus centrum\indexS{centrum librationis} librationis 
\textit{B}, et brachium\indexS{brachium librae} oppositum \textit{BE}, 
sit oneratum pondere \textit{E}. Quaeritur situs aequilibrii,\indexS{aequilibrium}\indexS{situs aequilibrii}
seu maximus descensus\indexS{descensus centri gravitatis maximus}
centri gravitatis communis.\indexS{centrum gravitatis commune} %
\pend
\pstart
Per \textit{B} ducatur recta
\edtext{verticalis \textit{AB}, cui occurrat}{%
\lemma{verticalis}%
\Bfootnote{%
\textit{(1)}~, cui occu %
\textit{(2)}~\textit{AB}, cui occurrat~\textit{L}}} 
planum\indexS{planum inclinatum} \textit{AD}
\edtext{in \textit{A}. Datus}{%
\lemma{in \textit{A}}%
\Bfootnote{%
\textit{(1)}~et ex punctis \textit{C} %
\textit{(2)}~. Datus~\textit{L}}} 
est angulus \textit{BAD}, 
\edtext{praeterea anguli}{%
\lemma{praeterea}%
\Bfootnote{%
\textit{(1)}~angulus \textit{CD}
\textit{(2)}~anguli~\textit{L}}}
\textit{ADC}, et \textit{BFC} sunt 
\edtext{recti. Caeterum}{%
\lemma{recti. Caeterum}%
\Bfootnote{%
\textit{(1)}~Ex
\textit{(2)}~Caeterum~\textit{L}}}
ex punctis \textit{C}, \textit{G}, \textit{D} agantur in ipsam \textit{AB} normales \textit{CH}, \textit{GL}, \textit{DM}. 
\edtext{Et quia}{%
\lemma{Et quia}%
\Bfootnote{%
\textit{(1)}~qua
\textit{(2)}~quia~\textit{L}}}
dantur rectae \textit{CG}, \textit{CD}, \textit{GD} 
daturque etiam angulus \textit{HCD}, vel \textit{LGD} vel \textit{MDG}, 
ideo dabuntur rectae \textit{HL}, \textit{HM}, \textit{LM}. Sed \textit{AH}, adeoque 
\edtext{\textit{AL}, et \textit{AM}}{%
\lemma{\textit{AL}, et}%
\Bfootnote{%
\textit{(1)}~\textit{M}
\textit{(2)}~\textit{AM}~\textit{L}}}
variantur, pro altiore aut depressiore globi situ. Assumto igitur situ globi aliquo, eo scilicet qui quaeritur eo
\edtext{ipso habetur}{%
\lemma{ipso}%
\Bfootnote{%
\textit{(1)}~fit et p %
\textit{(2)}~habetur~\textit{L}}} 
et punctum \textit{F}, ducendo scilicet ex \textit{B} rectam \textit{BF} quae globum tangat; adeoque habetur et punctum \textit{E}, ex
\edtext{quo agatur}{%
\lemma{quo}%
\Bfootnote{%
\textit{(1)}~demittatur %
\textit{(2)}~agatur~\textit{L}}} 
\textit{EN} normalis ad \textit{AB}, et recta \textit{LN} secetur in puncto \textit{P}, sic ut sit \textit{NP} ad \textit{LP} ut globus \textit{C} ad
\edtext{pondus \textit{E}. Itaque}{%
\lemma{pondus \textit{E}.}%
\Bfootnote{%
\textit{(1)}~Et punctum \textit{P} respondet c %
\textit{(2)}~Itaque~\textit{L}}} 
cum punctum \textit{N} responderit centro gravitatis\indexS{centrum gravitatis} ipsius ponderis \textit{E}, et punctum \textit{L} responderit
\edtext{ipsi \textit{G}}{%
\lemma{}%
\Bfootnote{ipsi \textit{G} \textit{erg.~L}}} 
centro gravitatis\indexS{centrum gravitatis}
\edtext{ipsius Globi \textit{C}.}{%
\lemma{ipsius}%
\Bfootnote{%
\textit{(1)}~ponderis %
\textit{(2)}~Globi \textit{C}.~\textit{L}}} 
Patet \textit{P} respondere centro\indexS{centrum gravitatis commune}
\edtext{gravitatis communi}{%
\lemma{gravitatis}%
\Bfootnote{%
\textit{(1)}~communis %
\textit{(2)}~communi~\textit{L}}}
ejusque altitudinem\indexS{altitudo centri gravitatis} vel depressionem\indexS{depressio centri gravitatis}
\edtext{exprimere. Quaeritur}{%
\lemma{exprimere.}%
\Bfootnote{%
\textit{(1)}~Debet igitur %
\textit{(2)}~Quaeritur~\textit{L}}} 
ergo situs globi \textit{C} adeoque ponderis \textit{E}
\edtext{talis, \bracks{ut} \textit{P}}{%
\lemma{talis,}%
\Bfootnote{%
\textit{(1)}~ut recta %
\textit{(2)}~\bars{ut \textit{erg. Hrsg.}} \textit{P}~\textit{L}}} 
prodeat omnium\indexS{situs centri gravitatis omnium possibilium infimus}
possibilium infimum; seu ut recta \textit{AP} sit omnium possibilium minima, posito \textit{P} cadere supra \textit{A}, vel omnium possibilium maxima, posito \textit{P} cadere infra \textit{A}.
\pend
\pstart
Superest ergo 
\edtext{ut ipsius}{%
\lemma{ut}%
\Bfootnote{%
\textit{(1)}~ipsa
\textit{(2)}~ipsius~\textit{L}}}
rectae \textit{AP} 
magnitudinem\indexS{magnitudo Analytice expressa} exprimamus Analytice,
ac deinde
\edtext{per nostram Methodum de Maximis et Minimis,\indexS{Methodus de Maximis et Minimis}}{%
\applabel{FnC:35_10_14_005_Maximis}%
\lemma{per \bracks{...} Minimis}%
\Cfootnote{Siehe \textsc{G.\,W.~Leibniz}, \cite{02077}\textit{Nova methodus pro maximis et minimis}, \cite{01023}\textit{AE} Oktober 1684, S.\;467--473.}}
situm debitum hinc eruamus, adeoque relationem quoque inter pondera %
\edlabel{35_10_14_005_1a}%
\edtext{}{%
\xxref{35_10_14_005_1a}{35_10_14_005_1b}%
\lemma{situsque.}%
\Bfootnote{%
\bracks{/}
\textit{(1)}~ \textit{AH} vocetur \textit{y} %
\textit{(2)}~\textit{AM} vocetur \textit{y}. Jam ob datum angulum %
\textit{(3)}~Ex puncto \textit{C} demittatur %
\textit{(a)}~recta
\textit{(b)}~perpendicularis
\textit{(c)}~perpendiculariter~\textit{L}}}%
situsque.
\pend
\pstart
Ex puncto \textit{C} demittatur perpendiculariter\edlabel{35_10_14_005_1b}
\textit{CQ} in ipsam \textit{MD}. Et globi radius \textit{CD}, vocetur \textit{r}, ipsa 
\edtext{\textit{CQ} vocetur \textit{c} (aequalis}{%
\lemma{\textit{CQ} vocetur}%
\Bfootnote{%
\textit{(1)}~\textit{q}  %
\textit{(2)}~\textit{c} %
\textit{(a)}~et ipsa \textit{QD} v %
\textit{(b)}~(aequalis~\textit{L}}}
ipsi \textit{HM}) et recta \textit{QD} vocetur \textit{q}, quae omnes sunt datae. Ipsa autem \textit{AM} quaesita vocetur \textit{y}, 
et ob triangula similia \textit{AMD}, et \textit{DQC} erit \textit{DM} ad \textit{AM} seu \textit{y}, ut \textit{CQ} seu \textit{c} ad \textit{QD} seu \textit{q}. Adeoque \textit{DM} erit $yc:q$ et \textit{AH} erit $y+c$ et \textit{CH} erit $\overline{yc:q}-q$. Porro quia ex puncto \textit{B} tangens
\edtext{\textit{BF}}{%
\lemma{}%
\Bfootnote{\textit{BF} \textit{erg.~L}}} 
duci debet ad circulum \textit{FD} centro \textit{C} descriptum ideo situs ipsius \textit{BF} debet esse talis, ut normalis \textit{CF} ex \textit{C} in \textit{BF} acta aequetur radio circuli seu ipsi \textit{CD}. Ex \textit{F} agatur in \textit{AB}, normalis \textit{FS},
\edtext{et in \textit{CH} normalis \textit{FR}\bracks{,}}{%
\lemma{}%
\Bfootnote{et in \textit{CH} normalis \textit{FR} \textit{erg.~L}}}
et \textit{CH} producatur dum ipsi \textit{FE} occurrat
\edtext{in \textit{T}. Et triangula}{%
\lemma{in \textit{T}.}%
\Bfootnote{%
\textit{(1)}~Jam ob Triangula simi %
\textit{(2)}~Et triangula~\textit{L}}} 
\textit{FSB} et \textit{TFC} sunt similia.
Item triangula \textit{TFC} et \textit{FRC} sunt similia, itaque triangula \textit{FSB} et \textit{FRC} sunt simila. Seu $FS:SB:BF\squaredots FR:RC:CF$. %
\pend
\pstart
Hinc erit 
\rule[0ex]{0pt}{4ex}%
\edtext{$SB:SF:\overset{BF}{\sqrt{SB^2+SF^2}} \squaredots \overline{SB+BH}$}{%
\lemma{}%
\Bfootnote{$SB:SF:\sqrt{SB^2+SF^2}$ \bars{\textit{BF}~erg.} $\squaredots\;\overline{SB+BH}$~\textit{L}}}
(seu \textit{FR}) $:\overline{CH-SF}$, (seu \textit{RC}) $:CF$ seu 
\edlabel{35_10_14_005_2a}%
\edtext{}{\xxref{35_10_14_005_2a}{35_10_14_005_2b}%
\lemma{}%
\Bfootnote{%
$\sqrt{\,\overline{SB+BH}^2+\overline{CH-SF}^2}$.
\bars{Habemus ergo hinc: $SB:SF\squaredots \overline{SB+BH}:CH-SF$ seu $SB\cdot CH-SB\cdot SF=SB\cdot SF+SF\cdot BH$ seu $2SB\cdot SF=SB\cdot CH-SF\cdot BH$ rursus $SB:$
\textit{(1)}~$SB^2:\overline{SB}^2$ %
\textit{(2)}~\bars{$\overline{SB^2+SF^2}$ \textit{ändert Hrsg.}} $\squaredots\;\overline{SB:BH}:CF$. \textit{gestr. u.~wieder gültig gemacht}}
\textit{(1)}~\bars{Habemus ergo 2 \textit{streicht Hrsg.}} %
\textit{(2)}~Rursus %
\textit{(3)}~Rursus %
\textit{(3)}~$CF^2$ %
\textit{(4)}~$SB^2+2SB\cdot BH+BH^2+CH^2-2CH\cdot EF+SF^2+$%
\textit{(5)}~rursus jungatur \textit{BC} quae est data, ut et \textit{FC}, fiet $BF=\sqrt{BC^2+CF^2}$. 
Datur ergo $BF^2=SB^2+SF^2=BC^2+CF^2$. Fiet ergo $SB^2+SF^2+2SB=BC^2+CF^2+S$ %
\textit{(6)}~Rursus~\textit{L}}}%
\rule[-.5ex]{0pt}{4ex}%
$\sqrt{\,\overline{SB+BH}^2+\overline{CH-SF}^2}$. Habemus ergo hinc: $SB:SF\squaredots \overline{SB+BH}:CH-SF$ seu 
$SB\cdot CH-SB\cdot SF=SB\cdot SF+SF\cdot BH$ seu $2SB\cdot SF=SB\cdot CH-SF\cdot BH$ rursus $SB:\bracks{\sqrt{\,\overline{SB^2+SF^2}}} \squaredots \overline{SB:BH}:CF$. Rursus%
\edlabel{35_10_14_005_2b}
jungatur \textit{BC} quae est data ut et \textit{FC}, fiet 
\edtext{$BF^2$ seu $SB^2+SF^2=BC^2+\bracks{CF^2}$, habetur}{%
\lemma{$BF^2$}%
\Bfootnote{%
\textit{(1)}~(seu
\textit{(2)}~seu $SB^2+SF^2=BC^2\,+$ \bars{$CF^2)$, \textit{ändert Hrsg.}} habetur~\textit{L}}}
ergo \textit{BF}.
\pend
\pstart
\href{\Scan{LH_35_10_14_005v}}%
{\lbrack5~v\textsuperscript{o}\rbrack}\ 
Jam $SB:BF\squaredots \overline{SB+BH}:CF$. 
Fit $SB\cdot CF=SB\cdot BF+BF\cdot BH$. Ergo fiet $SB=BF\cdot BH:\overline{CF-BF}$ seu $SB=\sqrt{BC^2+CF^2}\cdot BH:CF-\sqrt{BC^2+CF^2}$, 
ubi \textit{BC} intendit incognitam, nam $BC^2=BH^2+CH^2$,
et $CH=
\edtext{\overline{yc:q}-q$. Porro datur \textit{BE} quam vocabimus \textit{e}. Fiet ergo ut \textit{BE} seu \textit{e}, ad \textit{BF} ita \textit{BN} ad \textit{SB}. Ergo $BN=SB\cdot e:BF$ seu $BN=BH.e:\overline{CF-BF}$. Porro $BH=AH-AB$}{%
\lemma{}%
\Bfootnote{%
$\overline{yc:q}-q$. \bars{Ipsa autem \textit{BH} datur quam vocabimus \textit{h}.~\textit{gestr.}}
Porro datur 
\textit{(1)}~\textit{CE} quam vocabimus \textit{Ce}. Fiet ergo ut \textit{BE} seu \textit{Ce}, ad \bracks{...} Ergo  $BN=SB\cdot Ce:BF$ %
\textit{(2)}~\textit{BE} \bracks{...} $BN=SB\cdot e:BF$
\bars{seu $BN=BH.e:\overline{CF-BF}$.~\textit{erg.}}
\textit{(a)}~Jam datur \textit{BL} seu \textit{BH} vel \textit{h}, \textit{-HL} seu \textit{l}. Ergo $BL=h-l$ %
\textit{(b)}~Porro
\textit{(aa)}~\textit{CH}
\textit{(bb)}~$BH=AH-AB$~\textit{L}}}
et $AH=AM+MH=y+c$. Sit $AB=b$, et fiet $BH=y+c-b$. Porro datur \textit{HL}, quam vocemus \textit{l}, fiet $BL=y+c-b-l$
\edtext{et $AL=y+c-l$}{%
\lemma{}%
\Bfootnote{et $AL=y+c-l$ \textit{erg.~L}}} 
et fiet \rule[-.5ex]{0pt}{4ex}%
$NL=BN-BL=\overline{y+c-b}\;e:\overline{CF(-\sqrt{BC^2+CF^2}\text{\ seu})-\sqrt{BH^2+CH^2+CF^2}} -\overline{y+c-b-l}$
seu \rule[-.5ex]{0pt}{5ex}%
$NL=\overline{y+c-b}\;e:\overline{r-\sqrt{\,\overline{y+c-b}^2+\overline{\overline{yc:q}-q}^2+rr}}-\overline{y+c-b-l}$. 
\edtext{Jam \textit{LP}}{%
\lemma{Jam \textit{LP}}%
\Bfootnote{%
\textit{(1)}~\textit{LH}
\textit{(2)}~\textit{LH}
\textit{(3)}~\textit{LP}~\textit{L}}}
est ad \textit{NL} ut pondus in \textit{E} quod vocabimus \textit{p} est 
\edtext{ad pondus in \textit{E}}{\lemma{}\Bfootnote{ad pondus \bars{\textit{e} \textit{gestr.}} in \textit{E}~\textit{L}}}
et pondus globi quod vocabimus \textit{y} simul, seu \textit{LP} ad 
\edtext{\textit{NL} ut \textit{p} ad}{%
\lemma{\textit{NL} ut}%
\Bfootnote{%
\textit{(1)}~\textit{c} ad %
\textit{(2)}~\textit{p} ad~\textit{L}}} 
$\overline{p+g}$, adeoque tandem fiet %
\pend
\pstart
\noindent
\rule[-2mm]{0pt}{10mm}%
\edtext{%
$LP=\overline{\overline{\overline{y+c-b}\;%
e:\overline{r-\sqrt{\,\overline{y+c-b}^2+\overline{\overline{yc:q}-q}^2+rr}}}-\overline{y+c-b-l}}\;p:\overline{p+g}$.%
}{%
\lemma{$e$}%
\Bfootnote{%
\textit{(1)}~$p:\overline{p+g}$
\textit{(2)}~\rule[-.5ex]{0pt}{4ex}%
$:\overline{r-\sqrt{\,\overline{y+c-b}^2}}$~\textit{L}}} 
\pend
\pstart
\noindent
\edtext{Et $AP=$}{%
\lemma{$p:\overline{p+g}$}%
\Bfootnote{%
\textit{(1)}~$LP=$ %
\textit{(2)}~Et $AP=$~\textit{L}}}
$\overline{\overline{\overline{\overline{y+c-b}\;e:\overline{r-\sqrt{\,\overline{y+c-b}^2+\overline{\overline{yc:q}-q}^2+rr}}}-\overline{y+c-b-l}}\;p:\overline{p+g}}+y+c-l$. %
\pend
\pstart
\noindent
Quae \textit{AP} debet esse omnium possibilium minima vel
\edtext{maxima. Quoniam}{%
\lemma{maxima}%
\Bfootnote{%
\textit{(1)}~et fit $0=d$ %
\textit{(2)}~. Quoniam~\textit{L}}} 
autem compendii causa multa possumus assumere pro arbitrio ut magnitudinem ipsius \textit{AB} seu rectae \textit{b}, item ipsius \textit{BE}, seu rectae \textit{e}, item ipsius radii \textit{r} hinc formula\indexS{formula irrationalis} irrationalis reddi poterit extrahibilis. %
\pend
\pstart
Generaliter pro $y+c$ substituatur $v=AH$, et pro $yc:q$ ponatur $\overline{v-c}\;c:q$. Ita formula extrahenda erit $v^2-2vb+bb+\overline{v^2-2vc+cc}\;c^2:q^2-2\overline{v-c}\;c+qq+rr$. Seu $vv+\overline{-2b-\overline{2c^3:qq}-2c}:\overline{1+\overline{cc:qq}}\;v+\overline{bb+\overline{c^2:qq}+2cc+qq+rr}:\overline{1+\overline{cc:qq}}$.%
\rule[-.5ex]{0pt}{3.5ex}
Quae ut fiat extrahibilis si comparemus cum $vv-2Kv+KK$ debet esse $\overline{b+\overline{c^3:qq}+c:\overline{1+\overline{cc:qq}}}^2=\overline{bb+\overline{c^4:qq}+2cc+qq+rr}:\overline{1+\overline{cc:qq}}$.%
\rule[-.5ex]{0pt}{3.5ex}
\pend
\begin{edfigure}%
\includegraphics[scale=1]{%
gesamttex/edit_VIII-4/images/LH_35_10_14_005_d1_005v.pdf}
\edcaption{\lbrack\textit{Fig.~1}\rbrack}
\end{edfigure}
